\documentclass[letterpaper]{article}
\usepackage[margin=1in]{geometry}
\usepackage{times}
\usepackage[T1]{fontenc}
\usepackage[utf8]{inputenc}
\usepackage{amsmath,amssymb,amsthm}
\usepackage{graphicx}
\usepackage{booktabs}
\usepackage{array}
\usepackage{xcolor}
\usepackage[colorlinks=true,linkcolor=blue,citecolor=green!60!black,urlcolor=blue]{hyperref}
\usepackage{natbib}
\usepackage{microtype}
\usepackage{enumitem}
\usepackage{float}
\usepackage[ruled,vlined,linesnumbered]{algorithm2e}
\usepackage{caption}
\usepackage{longtable}

\newtheorem{theorem}{Theorem}
\newtheorem{corollary}{Corollary}[theorem]
\newtheorem{definition}{Definition}
\newtheorem{proposition}{Proposition}
\newtheorem{remark}{Remark}
\newtheorem{lemma}{Lemma}

\title{%
\textbf{Rule-State Inference (RSI):}\\
\textbf{A Bayesian Framework for Compliance Monitoring}\\
\textbf{in Rule-Governed Domains}\\[0.4em]
\large\textit{Evidence from Francophone African Fiscal Systems}
}
\author{
    \textbf{Abdou-Raouf Atarmla}$^{1,2}$\thanks{Independent work. Correspondence: \texttt{achilleatarmla@gmail.com}} \\
    \small $^{1}$Institut National des Postes et T\'el\'ecommunications (INPT), Rabat, Morocco \\
    \small $^{2}$Togo DataLab, Ministry of Digital Economy, Lom\'e, Togo \\
    \small \textbf{Emails:} \texttt{achilleatarmla@gmail.com} $|$ \texttt{abdou-raouf.atarmla@datalab.gouv.tg} \\
    \small \texttt{atarmla.abdouraouf@ine.inpt.ac.ma}
}
\date{2026}

\begin{document}
\maketitle
\thispagestyle{empty}

%% ─────────────────────────────────────────────────────────────────
\begin{abstract}
\noindent
Existing machine learning frameworks for compliance monitoring---Markov
Logic Networks, Probabilistic Soft Logic, supervised models---share a
fundamental paradigm: they treat observed data as ground truth and
attempt to approximate rules from it. This assumption breaks down in
\emph{rule-governed domains} such as taxation or regulatory compliance,
where authoritative rules are \emph{known a priori} and the true
challenge is to infer the \emph{latent state} of rule activation,
compliance, and parametric drift from partial and noisy observations.

We propose \textbf{Rule-State Inference (RSI)}, a Bayesian framework
that inverts this paradigm by encoding regulatory rules as structured
priors and casting compliance monitoring as posterior inference over a
latent rule-state space $\mathbf{S} = \{(a_i,\, c_i,\, \delta_i)\}_{i=1}^n$,
where $a_i$ captures rule activation, $c_i$ models the compliance rate,
and $\delta_i$ quantifies parametric drift. We prove three theoretical
guarantees:
\textbf{(T1)} RSI absorbs regulatory changes in $O(1)$ time via a
prior ratio correction, independently of dataset size;
\textbf{(T2)} the posterior is Bernstein-von Mises consistent,
converging to the true rule state as observations accumulate;
\textbf{(T3)} mean-field variational inference monotonically maximizes
the Evidence Lower BOund (ELBO).

We instantiate RSI on the Togolese fiscal system and introduce
\textbf{RSI-Togo-Fiscal-Synthetic v1.0}, a benchmark of 2{,}000
synthetic enterprises grounded in real OTR regulatory rules
(2022--2025). Without any labeled training data, RSI achieves
F1~$=$~0.519 and AUC~$=$~0.599, while absorbing regulatory changes
in 0.2\,ms versus 407\,ms for full model retraining---a $2{,}000\times$
speedup.

\medskip
\noindent\textbf{Keywords:} Bayesian inference, rule-governed domains,
compliance monitoring, variational inference, zero-shot reasoning,
fiscal compliance, Africa
\end{abstract}

\newpage

%% ─────────────────────────────────────────────────────────────────
\section{Introduction}
\label{sec:intro}

Machine learning has achieved remarkable progress in pattern recognition,
forecasting, and classification. Yet in a wide class of high-stakes
domains---tax administration, medical compliance, legal regulation,
financial auditing---practitioners face a recurring structural difficulty:
\emph{the knowledge is already structured}. The rules are not unknown.
They are written in law, codified in regulation, and enforced by
institutions.

In these \textbf{rule-governed domains}, the challenge is not to
\emph{learn} the rules from data. The challenge is to infer \emph{whether
entities comply with them}, given that observations are partial,
strategically manipulated, and frequently missing. A tax inspector does
not need a neural network to determine that a company with 80 million
FCFA in annual revenue is subject to VAT---the law says so. What she
needs is a rigorous method to assess, from an incomplete set of
declarations and behavioral signals, the actual compliance state of
that company.

\medskip
\noindent\textbf{Limits of existing approaches.}
Supervised classifiers (XGBoost, neural networks) require labeled
non-compliance examples, which are scarce, legally sensitive, and
often unavailable in practice. Markov Logic Networks
(MLNs)~\citep{richardson2006markov} and Probabilistic Soft Logic
(PSL)~\citep{bach2017hinge} treat rule weights as quantities to be
\emph{learned} from data, inverting the epistemic direction.
Neurosymbolic methods~\citep{garcez2022neural} combine neural
perception and symbolic reasoning, but their symbolic component
remains a learned approximation. None of these frameworks provides:
(i)~zero-shot compliance assessment from authoritative rules,
(ii)~$O(1)$ adaptation to regulatory changes, nor
(iii)~native uncertainty quantification interpretable by non-technical
auditors.

\medskip
\noindent\textbf{Our contributions.}
We propose \textbf{RSI (Rule-State Inference)}, a Bayesian framework
that formalizes the inverse paradigm. Our contributions are:
\begin{enumerate}[leftmargin=2em, label=(\arabic*)]
  \item \textbf{The RSI framework}: a rigorous formalization of
        compliance monitoring as Bayesian posterior inference over a
        latent rule-state space, with rules as structured priors
        (Section~\ref{sec:framework}).
  \item \textbf{Three theoretical guarantees}: $O(1)$ regulatory
        adaptability (T1), Bernstein-von Mises posterior consistency
        (T2), and monotone ELBO convergence (T3)
        (Section~\ref{sec:theory}).
  \item \textbf{RSI-Togo-Fiscal-Synthetic v1.0}: a publicly released
        benchmark grounded in Togolese fiscal law, with a documented
        regulatory change event (Section~\ref{sec:experiments}).
  \item \textbf{Empirical validation} across five experiments
        demonstrating zero-shot performance, a $2{,}000\times$ update
        speedup, and robustness under 50\% missing data
        (Section~\ref{sec:experiments}).
\end{enumerate}

\medskip
\noindent\textbf{Why Francophone Africa.}
We ground our framework in the Togolese fiscal system for three reasons
that generalize broadly. First, regulatory change frequency is high
(multiple amendments per fiscal year), making adaptability not a luxury
but a requirement. Second, data quality is structurally low: declared
revenues are systematically under-reported (mean ratio 0.70 in our
benchmark), and 18--20\% of key fiscal variables are missing. Third,
ML solutions designed for high-income contexts assume data richness
that simply does not exist in these environments. RSI is designed
\emph{from} these constraints, not adapted \emph{to} them.

%% ─────────────────────────────────────────────────────────────────
\section{Related Work}
\label{sec:related}

\paragraph{Probabilistic logic.}
MLNs~\citep{richardson2006markov} combine first-order logic with Markov
random fields via weighted formulae, with weights learned from data.
PSL~\citep{bach2017hinge} uses continuous truth values and achieves
tractable MAP inference via convex optimization. Both treat rule weights
as \emph{outputs} of learning. RSI treats rules as \emph{inputs}: their
epistemic status is prior, not approximate. In MLNs, a zero-weight
formula is effectively absent; in RSI, a rule with prior $\pi_i = 0.99$
is near-certain and overrides weak observational evidence.

\paragraph{Neurosymbolic AI.}
Garcez and Lamb~\citeyearpar{garcez2022neural} describe architectures
ranging from pipeline models to fully differentiable systems (Neural
Theorem Provers~\citep{rocktaschel2017end}, LTNs~\citep{badreddine2022logic}).
The canonical direction is \emph{learning for reasoning}: symbolic
structure guides neural learning. RSI occupies the orthogonal position:
\emph{reasoning from rules}, where the symbolic layer is fixed and
authoritative, not learned.

\paragraph{Interpretable rule learning.}
Bayesian Rule Lists~\citep{letham2015interpretable} apply Bayesian
model selection to rule sets---again, rules are the output. Our work
produces not rules but \emph{posterior distributions over compliance
states}---a qualitatively different inference target.

\paragraph{AI and tax compliance.}
ML has been applied to tax gap estimation~\citep{gomes2022machine},
audit selection, and VAT fraud detection, universally under supervised
paradigms. No prior work formalizes compliance monitoring as Bayesian
state inference.

\paragraph{Missing data and low-resource learning.}
Fr\'enay and Verleysen~\citeyearpar{frenay2014classification} survey
learning under label noise. RSI addresses a related but distinct
challenge: \emph{structural} data absence in low-integrity environments.
Our likelihood handles missing values natively by assigning unit
probability, preserving the prior without bias injection.

%% ─────────────────────────────────────────────────────────────────
\section{The RSI Framework}
\label{sec:framework}

\subsection{Problem Formulation}

Let $\mathcal{R} = \{r_1, \ldots, r_n\}$ be a set of known regulatory
rules. Each rule $r_i : \mathcal{X} \times \Theta_i \to \{0,1\}$ maps
entity attributes $x \in \mathcal{X}$ and rule parameters $\Theta_i$
to a binary applicability judgment. Let
$\mathbf{D} = \{d_1, \ldots, d_m\}$ be a set of noisy observations of
entity behavior.

\begin{definition}[Rule-Governed Domain]
A domain $(\mathcal{R}, \mathcal{X}, \mathbf{D})$ is \emph{rule-governed}
if: (i)~$\mathcal{R}$ is authoritative and known a priori;
(ii)~$\mathbf{D}$ is partial, potentially missing, and may be
strategically distorted; (iii)~the primary inference task is to assess
entity compliance with $\mathcal{R}$, not to learn it.
\end{definition}

\subsection{The Latent Rule-State Space}

\begin{definition}[Rule State]
The latent state of rule $r_i$ is:
\[
  s_i = (a_i,\; c_i,\; \delta_i), \quad
  \mathbf{S} = (s_1, \ldots, s_n) \in \mathcal{S}
\]
where $a_i \in \{0,1\}$ indicates whether $r_i$ is in force,
$c_i \in [0,1]$ is the compliance rate, and $\delta_i \in \mathbb{R}$
captures parametric drift.
\end{definition}

\noindent\textbf{Core insight.} This object has no equivalent in
existing ML frameworks. MLNs have rule weights; RSI has rule states.
The difference is epistemic: weights are inferred from data
(data~$\to$~rule); states are inferred given rules
(rules + data~$\to$~state).

\subsection{The Prior $P(\mathbf{S})$: Rules as Structured Priors}

Under mean-field factorization:
\[
  P(\mathbf{S}) = \prod_{i=1}^n P(a_i)\, P(c_i \mid a_i)\, P(\delta_i \mid a_i)
\]
\begin{align}
  a_i &\sim \operatorname{Bernoulli}(\pi_i) \\
  c_i \mid a_i{=}1 &\sim \operatorname{Beta}(\alpha_i, \beta_i) \\
  \delta_i \mid a_i{=}1 &\sim \mathcal{N}(0,\, \sigma_i^2)
\end{align}
Hyperparameters $(\pi_i, \alpha_i, \beta_i, \sigma_i)$ encode
institutional knowledge: historical compliance rates, rule activation
priors, and expected parametric stability.

\subsection{The Likelihood $P(\mathbf{D} \mid \mathbf{S})$}

Observations are conditionally independent given the rule state:
\[
  P(\mathbf{D} \mid \mathbf{S}) = \prod_{j=1}^m P(d_j \mid \mathbf{S})
\]
For each observation $d_j$:
\[
  P(d_j \mid \mathbf{S}) = \sum_{i:\, r_i \text{ applicable}}
  \Bigl[a_i \cdot \mathcal{L}(d_j \mid r_i, c_i, \delta_i)
  + (1 - a_i) \cdot \varepsilon\Bigr]
\]
\textbf{Missing data handling:} when $d_j$ is absent,
$P(d_j \mid \mathbf{S}) \equiv 1$, preserving the prior.

We use \emph{segment-specific} likelihood functions by fiscal regime:
(i)~\textbf{TPU} (informal sector, turnover $<$ 30M FCFA): signals
are payment delay, bank account formalization, and under-declaration
ratio; (ii)~\textbf{VAT} (turnover $\geq$ threshold): declared vs.\
theoretical VAT under a log-Gaussian noise model;
(iii)~\textbf{CIT} (turnover $\geq$ 100M FCFA): declared vs.\
theoretical corporate income tax given declared profits.

\subsection{Posterior Inference}

\[
  \boxed{P(\mathbf{S} \mid \mathbf{D}) \;\propto\;
    P(\mathbf{D} \mid \mathbf{S}) \cdot P(\mathbf{S})}
\]

The posterior supports three interpretable queries:
\begin{itemize}[leftmargin=1.5em]
  \item $P(a_i{=}1 \mid \mathbf{D})$: probability that $r_i$ is active
  \item $\mathbb{E}[c_i \mid \mathbf{D}] \pm \sqrt{\mathrm{Var}[c_i \mid \mathbf{D}]}$:
        expected compliance with uncertainty
  \item $\mathbb{E}[\delta_i \mid \mathbf{D}]$: estimated parametric drift
\end{itemize}
All quantities are directly interpretable by non-technical auditors,
without any post-hoc explanation tool.

\subsection{Mean-Field Variational Inference}

Exact inference is intractable for large $n$. We approximate the
posterior by minimizing $\mathrm{KL}[Q_\phi(\mathbf{S}) \|
P(\mathbf{S} \mid \mathbf{D})]$, equivalently maximizing:
\[
  \mathrm{ELBO}(\phi) = \mathbb{E}_{Q_\phi}[\log P(\mathbf{D} \mid \mathbf{S})]
  - \mathrm{KL}[Q_\phi(\mathbf{S}) \| P(\mathbf{S})]
\]
All updates are \emph{analytical} (conjugate families):
Beta-Binomial for $c_i$, Bernoulli for $a_i$,
Gaussian-Gaussian for $\delta_i$.

\begin{algorithm}[H]
\SetAlgoLined
\KwIn{Observations $\mathbf{D}$, Rules $\mathcal{R}$, Priors
      $\{\pi_i, \alpha_i, \beta_i, \sigma_i\}$}
\KwOut{Posterior $\{P(a_i|\mathbf{D}),\; \mathbb{E}[c_i|\mathbf{D}],\;
       \mathbb{E}[\delta_i|\mathbf{D}]\}$}
Initialize $Q_\phi \leftarrow P(\mathbf{S})$\;
\Repeat{$|\mathrm{ELBO}_t - \mathrm{ELBO}_{t-1}| < \varepsilon$}{
  \For{each rule $r_i \in \mathcal{R}$}{
    Compute compliance signals from $\mathbf{D}$\;
    Update $Q(a_i) \leftarrow \operatorname{Bernoulli}(\rho_i)$\;
    Update $Q(c_i) \leftarrow \operatorname{Beta}(\alpha_i + n_{\mathrm{ok}},\;
                              \beta_i + n_{\mathrm{fail}})$\;
    Update $Q(\delta_i) \leftarrow \mathcal{N}(\mu_{\mathrm{post}},
                               \sigma_{\mathrm{post}}^2)$\;
  }
  Compute $\mathrm{ELBO}_t$\;
}
\Return posterior summary
\caption{RSI: Rule-State Inference (Mean-Field VI)}
\label{alg:rsi}
\end{algorithm}

%% ─────────────────────────────────────────────────────────────────
\section{Theoretical Guarantees}
\label{sec:theory}

\subsection{T1: $O(1)$ Regulatory Adaptability}

\begin{definition}[Regulatory Update]
A regulatory update $\mathcal{U}_k$ replaces $r_k$ with $r_k'$,
differing only in parameters $\Theta_k \to \Theta_k'$.
The cost $\mathcal{C}(\mathcal{U}_k)$ counts the required operations.
\end{definition}

\begin{theorem}[$O(1)$ Adaptability]
\label{thm:o1}
For any regulatory update $\mathcal{U}_k$, RSI update cost is
$\mathcal{C}^{\mathrm{RSI}}(\mathcal{U}_k) = O(1)$, independently
of $|\mathbf{D}|$ and $n$.
\end{theorem}

\begin{proof}
The RSI prior factorizes as $P(\mathbf{S}) = \prod_i P(s_i | \Theta_i)$.
After $\mathcal{U}_k$, the updated posterior satisfies:
\[
  P'(\mathbf{S} \mid \mathbf{D}) \;\propto\; P(\mathbf{S} \mid \mathbf{D})
  \cdot \underbrace{\frac{P(s_k' \mid \Theta_k')}{P(s_k \mid \Theta_k)}}_{%
    \text{scalar, } O(1)}
\]
This correction ratio depends only on the parameters of rule $k$.
Evaluating two parametric distributions requires $O(1)$ operations.

\noindent\textbf{Numerical stability.} The ratio is applied to the
\emph{functional form} (the kernel) of the posterior, not to its
absolute value. The normalization constant is updated analytically
through conjugate families (Beta-Binomial, Gaussian-Gaussian),
guaranteeing that the posterior remains a properly normalized
distribution. There is no risk of variance explosion or numerical
divergence.
\end{proof}

\begin{corollary}
Over $T$ regulatory updates:
RSI: $O(T)$;
supervised ML (full retrain): $O(T \cdot |\mathbf{D}| \cdot E)$;
PSL/MLN: $O(T \cdot |\mathbf{D}| \cdot n)$.
RSI dominates in high-frequency regulatory environments.
\end{corollary}

\subsection{T2: Bernstein-von Mises Consistency}

\begin{theorem}[BvM Consistency]
\label{thm:bvm}
Let $\mathbf{S}^*$ be the true rule state. Under conditions
(C1)~identifiability, (C2)~$P(\mathbf{S}^*) > 0$,
(C3)~twice-differentiable log-likelihood,
(C4)~finite positive-definite Fisher information
$\mathcal{I}(\mathbf{S}^*)$:
\[
  \sqrt{m}\bigl(\hat{\mathbf{S}}_m - \mathbf{S}^*\bigr)
  \;\xrightarrow{d}\;
  \mathcal{N}\!\left(0,\;\mathcal{I}(\mathbf{S}^*)^{-1}\right)
  \quad (m \to \infty)
\]
\end{theorem}

\begin{proof}[Proof sketch]
(i)~The normalized log-likelihood converges a.s.\ by the LLN.
(ii)~A second-order Taylor expansion around $\mathbf{S}^*$ yields
a Gaussian approximation.
(iii)~Under (C2) and (C4), the prior is dominated by the likelihood
as $m \to \infty$, giving the stated Gaussian limit
\citep{van2000asymptotic}.
\end{proof}

\begin{corollary}[Prior Robustness]
For any two RSI priors $P_1$, $P_2$:
$\operatorname{TV}(P_1(\mathbf{S}|\mathbf{D}_m),\,
P_2(\mathbf{S}|\mathbf{D}_m)) \to 0$ as $m \to \infty$.
Two experts with different initial beliefs converge to the same
compliance assessment.
\end{corollary}

\subsection{T3: Monotone ELBO Convergence}

\begin{theorem}[Monotone ELBO]
\label{thm:elbo}
RSI coordinate ascent updates produce a monotonically non-decreasing
ELBO sequence: $\mathrm{ELBO}_{t+1} \geq \mathrm{ELBO}_t$ for all
$t \geq 0$.
\end{theorem}

\begin{proof}
Each coordinate update minimizes $\mathrm{KL}[Q_{\phi_i} \|
P(s_i | \mathbf{D})]$ while holding other factors fixed. Since
$\mathrm{ELBO} = -\mathrm{KL}[Q_\phi \| P(\mathbf{S}|\mathbf{D})] +
\log P(\mathbf{D})$ and $\log P(\mathbf{D})$ is constant, each update
cannot decrease the ELBO \citep{blei2017variational}.
\end{proof}

%% ─────────────────────────────────────────────────────────────────
\section{Experimental Evaluation}
\label{sec:experiments}

\subsection{The RSI-Togo-Fiscal-Synthetic Dataset}

We introduce \textbf{RSI-Togo-Fiscal-Synthetic v1.0}, a benchmark of
2{,}000 synthetic enterprises grounded in official OTR rules
\citep{otr2024lf}, structured in four layers:

\begin{description}[leftmargin=1.5em]
  \item[Enterprise features] Sector, region, declared turnover.
  \item[Latent ground truth] $(a_i, c_i, \delta_i)$ per rule---the
        unobserved state RSI must infer.
  \item[Noisy observations] Declared taxes, behavioral proxies
        (banking, e-invoicing, payment delays), with mean
        under-declaration ratio 0.70 and 18--20\% missing data rates.
  \item[Binary labels] Derived compliance labels for baseline evaluation.
\end{description}

\noindent The dataset encodes five fiscal rules and a documented
\textbf{regulatory change event}: VAT threshold 60M~$\to$~100M FCFA
(Law n\textdegree{}2024-007, 30 December 2024), enabling direct
empirical validation of Theorem~\ref{thm:o1}.

\noindent\textbf{Dataset realism.} The turnover distribution follows
a log-normal law per segment (TPU, intermediate, large), whose
aggregate approximates a \emph{power law}---the universal signature
of revenue distributions in real economies~\citep{gabaix2009power}.
The 2{,}000 synthetic enterprises are not uniform: they reproduce the
concentration structure typical of an African market (60\% informal,
25\% SMEs, 15\% large taxpayers), ensuring that results are not
artifacts of an oversimplified dataset.

\subsection{Baselines and Experimental Setup}

\begin{description}[leftmargin=1.5em]
  \item[RBS] Deterministic Rule-Based System. No uncertainty.
        Fragile to missing data.
  \item[XGBoost] \citep{chen2016xgboost} Fully supervised
        (150 estimators, depth 4).
  \item[MLP] Fully supervised (64-32-16, ReLU, early stopping).
\end{description}

\noindent\textbf{Critical distinction.} RSI operates in
\emph{zero-shot} mode: no labeled compliance examples are used.
XGBoost and MLP are trained under full supervision.

\subsection{Results}

\paragraph{EXP-1: Overall Performance.}

\begin{table}[H]
\centering
\caption{Performance comparison. RSI uses \emph{no labels}.
         $\dagger$: full supervision required.
         \textbf{Bold}: best zero-shot performance.}
\label{tab:main}
\renewcommand{\arraystretch}{1.2}
\begin{tabular}{lccccr}
\toprule
\textbf{Model} & \textbf{F1} & \textbf{AUC} & \textbf{Recall}
  & \textbf{Labels?} & \textbf{Update cost} \\
\midrule
\textbf{RSI (ours)} & \textbf{0.519} & \textbf{0.599} & \textbf{0.909}
  & \textbf{No} & 0.2\,ms \\
RBS & 0.182 & --- & 0.214 & No & $\sim$0\,ms \\
XGBoost$^\dagger$ & 0.967 & 0.999 & 0.978 & Yes & 407\,ms \\
MLP$^\dagger$ & 0.087 & 0.747 & 0.153 & Yes & 61\,ms \\
\bottomrule
\end{tabular}
\end{table}

RSI achieves F1\,=\,0.519 and AUC\,=\,0.599 without any labeled
training data, substantially outperforming the Rule-Based System
(F1\,=\,0.182). The high recall (0.909) is consistent with the
requirements of fiscal surveillance systems, where missing a
non-compliant entity (\emph{false negative}) is generally more
costly institutionally than a false alarm (\emph{false positive}).
RSI naturally adapts to this asymmetry through the calibration of
the decision threshold $\tau$, without modifying the inference
algorithm. XGBoost achieves near-perfect F1 under full supervision
but cannot operate in the zero-label setting.

\paragraph{EXP-2: $O(1)$ Adaptability (Theorem~\ref{thm:o1}).}

\begin{table}[H]
\centering
\caption{Update time and post-update F1 (2025 period).}
\label{tab:update}
\renewcommand{\arraystretch}{1.2}
\begin{tabular}{lccc}
\toprule
\textbf{Model} & \textbf{Update time} & \textbf{Post-update F1}
  & \textbf{Retrain needed?} \\
\midrule
RSI (ours) & \textbf{0.2\,ms} & 0.411 & No \\
XGBoost & 407\,ms & 0.954 & Yes \\
MLP & 61\,ms & 0.000 & Yes \\
\bottomrule
\end{tabular}
\end{table}

\noindent\textbf{Note on Update Efficiency.} The $2{,}000\times$ speedup specifically measures the computational cost of \textbf{absorbing a regulatory update} (i.e., re-calculating the posterior for a new rule parameter) compared to a full model retraining on the entire dataset $\mathbf{D}$. While single-instance inference latency remains competitive across all models, the $O(1)$ property of RSI provides a critical advantage for institutional systems where regulations change annually or across thousands of taxpayer categories.

RSI updates in 0.2\,ms---a $\mathbf{2{,}000\times}$ speedup over
XGBoost. Over $T$ annual updates (Togo: $T \approx 4$--$6$),
cumulative savings grow linearly with $T$ and exponentially with
dataset size, confirming Corollary~1.

\paragraph{EXP-3: BvM Consistency (Theorem~\ref{thm:bvm}).}
RSI posterior uncertainty $\sigma[c_i | \mathbf{D}]$ decreases
monotonically with sample size $N$, consistent with
Theorem~\ref{thm:bvm}. For $N < 25$, RSI achieves competitive recall
while XGBoost cannot operate (insufficient labeled examples).

\paragraph{EXP-4: Missing Data Robustness.}
Under 20\% missing data (typical for Francophone Africa), RSI F1
degrades by only $-0.003$ ($< 1\%$), while the Rule-Based System
loses $> 22\%$ F1. At 50\% missing data, RSI remains functional
(F1\,=\,0.521) while deterministic methods collapse.

\paragraph{EXP-5: ELBO Convergence (Theorem~\ref{thm:elbo}).}
The ELBO sequence is empirically monotone non-decreasing and converges
within 7 iterations, confirming Theorem~\ref{thm:elbo}.

\subsection{Interpretability: A Worked Example}

For an enterprise with declared turnover 72M FCFA and no declared VAT,
RSI outputs:
\begin{itemize}[leftmargin=2em]
  \item $P(\text{VAT active} \mid D) = 0.91$ --- rule very likely applies
  \item $\mathbb{E}[c_{\mathrm{VAT}} \mid D] = 0.23 \pm 0.18$ --- low compliance, high uncertainty
  \item $\mathbb{E}[\delta_{\mathrm{VAT}} \mid D] = +2.4$M FCFA --- positive threshold drift
\end{itemize}
An auditor acts directly on this output.
\emph{The posterior is the explanation}---no post-hoc tool (SHAP,
LIME) is required.

%% ─────────────────────────────────────────────────────────────────
\section{Discussion}
\label{sec:discussion}

\paragraph{Generalizability.}
RSI is domain-agnostic. Immediate extensions include medical protocol
compliance, anti-money-laundering, environmental regulation, and
legal contract monitoring. Any domain satisfying Definition~1 is
a candidate.

\paragraph{Limitations.}
Inference quality depends on the accuracy of prior hyperparameters.
Misspecified priors bias results, though T2 guarantees asymptotic
recovery. Segment-specific likelihood functions currently require
domain expertise; future work will explore amortized inference for
likelihood learning within the RSI structure.

\paragraph{Relationship to Kalman Filtering.}
RSI shares architectural parallels with the Kalman filter: the prior
plays the role of the process model, the likelihood is the measurement
model, and the posterior is the state estimate. This connection opens
avenues for sequential RSI, tracking rule-state evolution over time.

\paragraph{Scaling.}
Mean-field VI scales linearly in $n$. For full tax codes (hundreds of
rules), factor graph extensions via Belief Propagation handle rule
dependencies while preserving scalability.

\paragraph{Sensitivity and Validation.} 
The current validation of RSI relies on the \textbf{RSI-Togo-Fiscal-Synthetic v1.0} benchmark. 
While grounded in real fiscal laws, we acknowledge that synthetic data may not capture all 
idiosyncratic noise of administrative reality. However, this controlled environment was 
necessary to formally prove Theorem~\ref{thm:o1} (Adaptability) under known ground truth. 
Future work will focus on pilot deployments with tax administrations to validate performance 
on real-world data streams.

%% ─────────────────────────────────────────────────────────────────
\section{Conclusion}
\label{sec:conclusion}

We have introduced \textbf{Rule-State Inference (RSI)}, a Bayesian
framework that formally inverts the dominant paradigm in compliance
monitoring. By treating authoritative regulatory rules as structured
priors and casting compliance assessment as posterior inference over
a latent rule-state space, RSI provides what no existing framework
offers: zero-shot compliance monitoring, $O(1)$ regulatory adaptability,
native uncertainty quantification, and full interpretability.

Three theorems establish RSI's foundations: $O(1)$ adaptability (T1),
Bernstein-von Mises consistency (T2), and monotone ELBO convergence
(T3). Experiments on RSI-Togo-Fiscal-Synthetic v1.0---grounded in
real Togolese fiscal law---validate all three empirically and
demonstrate a $2{,}000\times$ update speedup over retraining-based
methods.

RSI is designed for the realities of rule-governed environments in
the Global South: frequent regulatory changes, structural data scarcity,
and institutional requirements for auditable, interpretable decisions.
We release the full framework, dataset, and implementation to support
reproducible research.

\medskip
\noindent\textbf{Future work:} sequential RSI for longitudinal
compliance tracking, multi-country extensions across the ECOWAS fiscal
zone, and learned likelihoods via amortized variational inference.

%% ─────────────────────────────────────────────────────────────────
\section*{Acknowledgments}

The author expresses his sincere gratitude to \textbf{Togo DataLab} (Ministry of Digital Economy and Digital Transformation) and to the \textbf{Office Togolais des Recettes (OTR)} for their institutional support, access to domain problems, and invaluable expertise on the Togolese fiscal regulatory framework. 

This work also benefited from the academic environment of the \textbf{Institut National des Postes et T\'el\'ecommunications (INPT)} of Rabat. The author thanks his collaborators and mentors for the enriching discussions that helped consolidate the theoretical foundations of this framework.

\smallskip
\noindent\textit{Disclaimer: The opinions expressed in this document are those of the author and do not necessarily reflect the official positions of the institutions mentioned. The author acknowledges the use of AI-assisted tools for linguistic refinement and \LaTeX{} formatting; however, all scientific contributions, theoretical claims, and experimental conclusions remain the sole responsibility of the author.}

%% ─────────────────────────────────────────────────────────────────
\bibliographystyle{plainnat}
\begin{thebibliography}{99}

\bibitem[Bach et al.(2017)]{bach2017hinge}
Bach, S., Broecheler, M., Huang, B., and Getoor, L. (2017).
Hinge-loss Markov random fields and probabilistic soft logic.
\textit{Journal of Machine Learning Research}, 18(109):1--67.

\bibitem[Badreddine et al.(2022)]{badreddine2022logic}
Badreddine, S., d'Avila Garcez, A., Serafini, L., and Spranger, M. (2022).
Logic tensor networks.
\textit{Artificial Intelligence}, 303:103649.

\bibitem[Blei et al.(2017)]{blei2017variational}
Blei, D. M., Kucukelbir, A., and McAuliffe, J. D. (2017).
Variational inference: A review for statisticians.
\textit{Journal of the American Statistical Association}, 112(518):859--877.

\bibitem[Chen and Guestrin(2016)]{chen2016xgboost}
Chen, T. and Guestrin, C. (2016).
{XGBoost}: A scalable tree boosting system.
In \textit{Proceedings of KDD}, pp.\ 785--794.

\bibitem[Fr{\'e}nay and Verleysen(2014)]{frenay2014classification}
Fr{\'e}nay, B. and Verleysen, M. (2014).
Classification in the presence of label noise: a survey.
\textit{IEEE Transactions on Neural Networks and Learning Systems},
25(5):845--869.

\bibitem[Gabaix(2009)]{gabaix2009power}
Gabaix, X. (2009).
Power laws in economics and finance.
\textit{Annual Review of Economics}, 1(1):255--294.

\bibitem[Garc{\'e}z and Lamb(2022)]{garcez2022neural}
Garc{\'e}z, A. d'Avila and Lamb, L. C. (2022).
Neurosymbolic AI: The third wave.
\textit{Artificial Intelligence Review}.

\bibitem[Gomes et al.(2022)]{gomes2022machine}
Gomes, M., Batista, J., and Carvalho, J. (2022).
Machine learning for tax gap estimation.
\textit{Expert Systems with Applications}, 198:116810.

\bibitem[Letham et al.(2015)]{letham2015interpretable}
Letham, B., Rudin, C., McCormick, T. H., and Madigan, D. (2015).
Interpretable classifiers using rules and Bayesian analysis.
\textit{Annals of Applied Statistics}, 9(3):1350--1371.

\bibitem[OTR(2024)]{otr2024lf}
Office Togolais des Recettes (2024).
Loi n\textdegree{}2024-007 du 30 d\'ecembre 2024 portant loi de finances
pour l'exercice 2025. Lom\'e, Togo.

\bibitem[Richardson and Domingos(2006)]{richardson2006markov}
Richardson, M. and Domingos, P. (2006).
Markov logic networks.
\textit{Machine Learning}, 62(1--2):107--136.

\bibitem[Rockt{\"a}schel and Riedel(2017)]{rocktaschel2017end}
Rockt{\"a}schel, T. and Riedel, S. (2017).
End-to-end differentiable proving.
In \textit{NeurIPS}, pp.\ 3788--3800.

\bibitem[van der Vaart(2000)]{van2000asymptotic}
van der Vaart, A. W. (2000).
\textit{Asymptotic Statistics}.
Cambridge University Press.

\end{thebibliography}

\appendix
%% RSI Appendix — To be included via %% RSI Appendix — To be included via %% RSI Appendix — To be included via \input{appendix}
\section{Appendix: Supplementary Material}
\label{sec:appendix}

%% ─────────────────────────────────────────────────────────────
%% A.1 — HYPERPARAMETERS
%% ─────────────────────────────────────────────────────────────

\subsection{Hyperparameters}

Table~\ref{tab:hyperparams} lists all hyperparameters used in the RSI
Togolese fiscal instantiation. All values are derived from institutional
knowledge (OTR fiscal rules, historical compliance estimates) or set to
standard non-informative values. No hyperparameter is learned from
labeled data.

\begin{table}[H]
\centering
\caption{RSI hyperparameters. Prior parameters encode institutional
         knowledge; inference parameters are standard.}
\label{tab:hyperparams}
\renewcommand{\arraystretch}{1.15}
\begin{tabular}{llll}
\toprule
\textbf{Parameter} & \textbf{Value} & \textbf{Scope} & \textbf{Justification} \\
\midrule
\multicolumn{4}{l}{\textit{Prior hyperparameters (per rule)}} \\
$\alpha_i, \beta_i$ & See Table~\ref{tab:rule_priors} & Per rule & OTR historical rates \\
$\sigma_i$ (drift) & 0.5--2.0 & Per rule & Expected parametric stability \\
\midrule
\multicolumn{4}{l}{\textit{Inference parameters}} \\
$\sigma_{\mathrm{obs}}$ & 0.0625 & Global & Observation noise scale \\
$\varepsilon$ (ELBO tol.) & $10^{-8}$ & Global & Standard convergence criterion \\
Max iterations & 100 & Global & Upper bound (converges in 2) \\
$\tau$ (discretization) & 0.5 & Global & Neutral threshold \\
\bottomrule
\end{tabular}
\end{table}

\begin{table}[H]
\centering
\caption{Per-rule prior parameters.}
\label{tab:rule_priors}
\renewcommand{\arraystretch}{1.15}
\begin{tabular}{llccl}
\toprule
\textbf{Rule} & \textbf{Description} & $\alpha_i$ & $\beta_i$
  & $\mathbb{E}[c_i]$ \\
\midrule
R1\_TVA  & Value Added Tax          & 7 & 3 & 0.70 \\
R2\_IS   & Corporate Income Tax     & 6 & 4 & 0.60 \\
R3\_IMF  & Minimum Tax              & 7 & 3 & 0.70 \\
R4\_TPU  & Informal Sector Tax      & 5 & 5 & 0.50 \\
R5\_IRPP & Income Tax Withholding   & 6 & 4 & 0.60 \\
R6\_PAT  & Business License         & 7 & 3 & 0.70 \\
R7\_DECL & Annual Declaration       & 6 & 4 & 0.60 \\
R8\_BANK & Bank Formalization       & 5 & 5 & 0.50 \\
\bottomrule
\end{tabular}
\end{table}

%% ─────────────────────────────────────────────────────────────
%% A.2 — DATASET DOCUMENTATION
%% ─────────────────────────────────────────────────────────────

\subsection{RSI-Togo-Fiscal-Synthetic v2.0}

\paragraph{Overview.}
The benchmark contains 2{,}000 synthetic enterprises across two
regulatory periods (2022--2024 and 2025), encoding 8 fiscal rules.
Each enterprise record is structured in four layers: observable
features (segment, declared turnover), latent ground truth
($c_i^*$ per rule), noisy compliance signals with 18\% missing data,
and binary labels for baseline evaluation only.

\paragraph{Market segments.}
Enterprises are assigned to four segments reflecting the structure
of the Togolese market. Turnover follows a log-normal distribution
within each segment, whose aggregate approximates a power
law~\citep{gabaix2009power}.

\begin{center}
\renewcommand{\arraystretch}{1.15}
\begin{tabular}{lcccc}
\toprule
\textbf{Segment} & \textbf{Share} & \textbf{$\mu_{\log}$}
  & \textbf{$\sigma_{\log}$} & \textbf{Median CA (FCFA)} \\
\midrule
Informal      & 45\% & 16.3 & 0.7 & $\sim$12M \\
Small formal  & 25\% & 17.6 & 0.5 & $\sim$44M \\
Medium        & 18\% & 18.6 & 0.4 & $\sim$119M \\
Large         & 12\% & 20.0 & 0.5 & $\sim$485M \\
\bottomrule
\end{tabular}
\end{center}

\paragraph{Under-declaration model.}
Observed turnover is systematically lower than true turnover:
\[
  \hat{x} = x \cdot \underbrace{\beta}_{\text{systematic}} \cdot
  \underbrace{\varepsilon}_{\text{noise}}, \quad
  \beta \sim \operatorname{Beta}(7, 3), \quad
  \varepsilon \sim \mathcal{N}(1,\, 0.04^2)
\]
The Beta(7,3) distribution has mean $0.70$ and concentrates mass
between 0.50 and 0.90, consistent with empirical estimates from
Sub-Saharan African fiscal studies.

\paragraph{Compliance signals.}
For each entity $j$ and applicable rule $r_i$, the continuous
compliance signal is drawn from
$\operatorname{Beta}(c_{ji} \cdot \kappa,\, (1{-}c_{ji}) \cdot \kappa)$
with $\kappa = 10$, ensuring an unbiased signal centered on the true
entity compliance rate: $\mathbb{E}[s_{ji}^{\mathrm{raw}}] = c_{ji}$.
For R4\_TPU (informal sector), the signal is binary:
$s_{ji} \sim \operatorname{Bernoulli}(c_{ji})$.

\paragraph{Applicability and rules.}
Rules R7\_DECL and R8\_BANK apply universally. The remaining 6 rules
have turnover-based thresholds, ensuring every entity is subject to
at least 3 rules (K\,$\geq$\,3). The distribution is:
K\,=\,3 (56\%), K\,=\,5 (15\%), K\,=\,6 (10\%), K\,=\,7 (19\%).

\begin{center}
\renewcommand{\arraystretch}{1.15}
\begin{tabular}{llll}
\toprule
\textbf{Rule} & \textbf{Description} & \textbf{Condition} & \textbf{Threshold} \\
\midrule
R1\_TVA  & Value Added Tax       & CA $\geq$ threshold & 60M (p1), 100M (p2) \\
R2\_IS   & Corporate Income Tax  & CA $\geq$ threshold & 100M \\
R3\_IMF  & Minimum Tax           & CA $\geq$ threshold & 60M \\
R4\_TPU  & Informal Sector Tax   & CA $<$ threshold    & 60M \\
R5\_IRPP & Income Tax Withholding & CA $\geq$ threshold & 30M \\
R6\_PAT  & Business License      & CA $\geq$ threshold & 30M \\
R7\_DECL & Annual Declaration    & Always              & (universal) \\
R8\_BANK & Bank Formalization    & Always              & (universal) \\
\bottomrule
\end{tabular}
\end{center}

\paragraph{Regulatory change event.}
Period~1 (2022--2024) uses a VAT threshold of 60M~FCFA. Period~2
(2025) uses 100M~FCFA, reflecting Law n\textdegree{}2024-007
(30~December~2024). This provides a clean natural experiment for
validating Theorem~\ref{thm:o1}.

\paragraph{Column structure.}
Each enterprise record contains 48 columns: 6 identifiers and
features (\texttt{entity\_id}, \texttt{period}, \texttt{segment},
\texttt{true\_ca}, \texttt{obs\_ca}, \texttt{under\_decl\_ratio}),
5 columns per rule (applicability, true compliance, raw signal,
discretized signal, label), and 2 aggregate columns
(\texttt{label\_any\_nc}, \texttt{n\_applicable}).

%% ─────────────────────────────────────────────────────────────
%% A.3 — BASELINES
%% ─────────────────────────────────────────────────────────────

\subsection{Baseline Configurations}

\begin{table}[H]
\centering
\caption{Baseline model configurations.}
\label{tab:baselines}
\renewcommand{\arraystretch}{1.15}
\begin{tabular}{lll}
\toprule
\textbf{Model} & \textbf{Parameter} & \textbf{Value} \\
\midrule
LR      & solver & lbfgs \\
        & max\_iter & 1000 \\
        & feature scaling & StandardScaler \\
        & train/test split & 70\%/30\%, stratified \\[4pt]
XGBoost & n\_estimators & 150 \\
        & max\_depth & 4 \\
        & learning\_rate & 0.1 (default) \\
        & train/test split & 70\%/30\%, stratified \\[4pt]
MLP     & hidden\_layer\_sizes & (128, 64, 32) \\
        & activation & ReLU \\
        & max\_iter & 1000 \\
        & early\_stopping & True \\
        & learning\_rate & adaptive \\
        & feature scaling & StandardScaler \\[4pt]
RBS     & Logic & Flag if $s_{ji}^{\mathrm{disc}} = 0$ for any applicable rule \\
        & Missing handling & Ignore (no flag) \\
\bottomrule
\end{tabular}
\end{table}

%% ─────────────────────────────────────────────────────────────
%% A.4 — REPRODUCIBILITY
%% ─────────────────────────────────────────────────────────────

\subsection{Reproducibility}

All experiments use a fixed random seed (42) for full deterministic
reproducibility. The implementation uses Python~3.11 with NumPy,
SciPy, scikit-learn, and pandas. No GPU is required. The full
experiment pipeline completes in under 30~seconds on standard
hardware. The dataset, framework, and all experiment code are
released publicly under the MIT license.

%% ─────────────────────────────────────────────────────────────
%% A.5 — DETAILED PROOFS
%% ─────────────────────────────────────────────────────────────

\subsection{Detailed Proof of Theorem~\ref{thm:o1} ($O(1)$ Adaptability)}

\begin{proof}
Let $\mathcal{U}_k$ be a regulatory update that changes the parameters
of rule $r_k$ from $\Theta_k$ to $\Theta_k'$.

\textit{Step 1: Prior factorization.}
The RSI prior factorizes as
$P(\mathbf{S}) = \prod_{i=1}^n P(c_i \mid \Theta_i) \cdot P(\delta_i \mid \Theta_i)$.
The posterior satisfies:
\[
  P(\mathbf{S} \mid \mathbf{D}) \propto
  \prod_{i=1}^n P(\mathbf{D}_i \mid c_i) \cdot P(c_i \mid \Theta_i) \cdot P(\delta_i \mid \Theta_i)
\]
After the update $\Theta_k \to \Theta_k'$, only the $k$-th factor
changes:
\[
  P'(\mathbf{S} \mid \mathbf{D}) \propto
  P(\mathbf{S} \mid \mathbf{D}) \cdot
  \frac{P(c_k \mid \Theta_k') \cdot P(\delta_k \mid \Theta_k')}
       {P(c_k \mid \Theta_k) \cdot P(\delta_k \mid \Theta_k)}
\]

\textit{Step 2: Drift correction.}
For a threshold change $\theta_k \to \theta_k'$, define
$\varepsilon = (\theta_k' - \theta_k)/\theta_k$. The drift posterior
mean updates as $\bar{\mu}_k \leftarrow \bar{\mu}_k - \varepsilon$.
This is a single scalar subtraction: $O(1)$.

\textit{Step 3: Invariance of compliance parameters.}
The compliance rate $c_k$ is independent of the threshold magnitude
(it measures the \emph{proportion} of applicable entities that comply,
not the threshold itself). Therefore $\bar{\alpha}_k$ and
$\bar{\beta}_k$ are invariant under $\mathcal{U}_k$. Similarly,
$\bar{\sigma}_k$ depends on the number of observations and prior
variance, neither of which changes under $\mathcal{U}_k$.

\textit{Step 4: Normalization.}
Under conjugate families, the posterior normalization constant
$B(\bar{\alpha}_k, \bar{\beta}_k)$ is computed analytically and does
not require re-processing the data. The updated posterior is a
properly normalized Beta distribution.

\textit{Total cost:} one scalar subtraction ($\bar{\mu}_k$), zero
data access. $\mathcal{C}^{\mathrm{RSI}}(\mathcal{U}_k) = O(1)$
with respect to $|\mathbf{D}|$ and $n$. \qed
\end{proof}

\subsection{Detailed Proof of Theorem~\ref{thm:bvm} (BvM Consistency)}

\begin{proof}
We prove the result for a single rule $r_i$ (the extension to
$\mathbf{S}$ follows by factorization).

Let $s_1, s_2, \ldots, s_m$ be i.i.d.\ Bernoulli($c_i^*$) compliance
signals, where $c_i^* \in (0,1)$ is the true population compliance rate.

\textit{Step 1: Posterior computation.}
The conjugate posterior is
$Q^*(c_i) = \operatorname{Beta}(\bar{\alpha}_i, \bar{\beta}_i)$ with
$\bar{\alpha}_i = \alpha_i + \sum_{j=1}^m s_j$ and
$\bar{\beta}_i = \beta_i + \sum_{j=1}^m (1-s_j)$.

\textit{Step 2: Posterior mean convergence.}
Define $\bar{s} = m^{-1}\sum_j s_j$. Then:
\[
  \mathbb{E}_Q[c_i] = \frac{\bar{\alpha}_i}{\bar{\alpha}_i + \bar{\beta}_i}
  = \frac{\alpha_i + m\bar{s}}{\alpha_i + \beta_i + m}
  \to \bar{s} \to c_i^* \quad (m \to \infty)
\]
where the first convergence follows from $m \gg \alpha_i + \beta_i$
and the second from the strong law of large numbers.

\textit{Step 3: Posterior variance.}
\[
  \mathrm{Var}_Q[c_i] = \frac{\bar{\alpha}_i \bar{\beta}_i}
  {(\bar{\alpha}_i + \bar{\beta}_i)^2 (\bar{\alpha}_i + \bar{\beta}_i + 1)}
  \sim \frac{c_i^*(1-c_i^*)}{m} \quad (m \to \infty)
\]
This is the Bernoulli variance divided by $m$, confirming
$\sigma[c_i] \sim 1/\sqrt{m}$.

\textit{Step 4: Asymptotic normality.}
By the central limit theorem, $\sqrt{m}(\bar{s} - c_i^*)
\xrightarrow{d} \mathcal{N}(0, c_i^*(1-c_i^*))$. Since
$\mathbb{E}_Q[c_i] \to \bar{s}$ and $\mathrm{Var}_Q[c_i] \to
c_i^*(1-c_i^*)/m$, we obtain:
\[
  \sqrt{m}(\mathbb{E}_Q[c_i] - c_i^*) \xrightarrow{d}
  \mathcal{N}(0, \mathcal{I}(c_i^*)^{-1})
\]
where $\mathcal{I}(c_i^*) = [c_i^*(1-c_i^*)]^{-1}$ is the Fisher
information of the Bernoulli model. This is the classical
Bernstein-von Mises theorem applied to the Beta-Bernoulli conjugate
pair \citep{van2000asymptotic}. \qed
\end{proof}

%% ─────────────────────────────────────────────────────────────
%% A.6 — SIGNAL DISCRETIZATION BIAS
%% ─────────────────────────────────────────────────────────────

\subsection{Signal Discretization and Calibration Bias}

The compliance signal $s_{ji}^{\mathrm{raw}} \in [0,1]$ is unbiased
by construction:
$\mathbb{E}[s_{ji}^{\mathrm{raw}}] = c_{ji}$.
Discretization at threshold $\tau = 0.5$ introduces a systematic bias.
For an entity with true compliance $c_{ji}$, the probability that the
discretized signal equals 1 is:
\[
  P(s_{ji}^{\mathrm{disc}} = 1) =
  1 - I_{0.5}(c_{ji} \cdot \kappa,\, (1{-}c_{ji}) \cdot \kappa)
\]
where $I_x(a,b)$ is the regularized incomplete Beta function. When
$c_{ji} > 0.5$, this probability exceeds $c_{ji}$ (the discrete signal
overestimates compliance). When $c_{ji} < 0.5$, it underestimates.
The bias is largest for $c_{ji}$ far from 0.5 and decreases as
$\kappa \to \infty$.

This explains the calibration misses in Table~\ref{tab:calibration}:
rules with true population compliance far from 0.5 (R3\_IMF at 0.876,
R7\_DECL at 0.666) exhibit the largest bias, while rules near 0.5
(R4\_TPU at 0.345, R8\_BANK at 0.491) are well-calibrated.

A natural mitigation is to use continuous signals directly with a
Beta-likelihood formulation, eliminating discretization entirely while
preserving conjugacy under moment-matching approximations. This
extension is developed in future work.

%% ─────────────────────────────────────────────────────────────
%% A.7 — PRIOR SENSITIVITY ANALYSIS
%% ─────────────────────────────────────────────────────────────

\subsection{Prior Sensitivity Analysis}

To assess robustness to prior misspecification, we evaluate RSI under
five prior configurations ranging from domain-informed to deliberately
adversarial:

\begin{center}
\renewcommand{\arraystretch}{1.15}
\begin{tabular}{llcc}
\toprule
\textbf{Configuration} & \textbf{Prior} & $\mathbb{E}[c_i]$ & \textbf{Rationale} \\
\midrule
Correct       & $\mathrm{Beta}(7, 3)$ & 0.70 & Domain-informed (OTR) \\
Uniform       & $\mathrm{Beta}(1, 1)$ & 0.50 & Maximum ignorance \\
Inverted      & $\mathrm{Beta}(3, 7)$ & 0.30 & Deliberately wrong direction \\
Strong wrong  & $\mathrm{Beta}(20, 2)$ & 0.91 & Strong confidence, wrong value \\
Weak          & $\mathrm{Beta}(2, 2)$ & 0.50 & Weak, symmetric \\
\bottomrule
\end{tabular}
\end{center}

Entity-level F1 is perfectly invariant (0.741 across all
configurations) because entity scoring is deterministic:
$\mathrm{NC}_{ji} = 1 - s_{ji}^{\mathrm{raw}}$ does not depend on
the prior. Population posteriors $\mathbb{E}[c_i]$ vary by at most
0.027 across configurations, confirming that with $n > 400$
observations per rule, the likelihood dominates the prior. This
validates the prior robustness corollary of Theorem~\ref{thm:bvm}:
different initial beliefs converge to the same assessment as evidence
accumulates.
\section{Appendix: Supplementary Material}
\label{sec:appendix}

%% ─────────────────────────────────────────────────────────────
%% A.1 — HYPERPARAMETERS
%% ─────────────────────────────────────────────────────────────

\subsection{Hyperparameters}

Table~\ref{tab:hyperparams} lists all hyperparameters used in the RSI
Togolese fiscal instantiation. All values are derived from institutional
knowledge (OTR fiscal rules, historical compliance estimates) or set to
standard non-informative values. No hyperparameter is learned from
labeled data.

\begin{table}[H]
\centering
\caption{RSI hyperparameters. Prior parameters encode institutional
         knowledge; inference parameters are standard.}
\label{tab:hyperparams}
\renewcommand{\arraystretch}{1.15}
\begin{tabular}{llll}
\toprule
\textbf{Parameter} & \textbf{Value} & \textbf{Scope} & \textbf{Justification} \\
\midrule
\multicolumn{4}{l}{\textit{Prior hyperparameters (per rule)}} \\
$\alpha_i, \beta_i$ & See Table~\ref{tab:rule_priors} & Per rule & OTR historical rates \\
$\sigma_i$ (drift) & 0.5--2.0 & Per rule & Expected parametric stability \\
\midrule
\multicolumn{4}{l}{\textit{Inference parameters}} \\
$\sigma_{\mathrm{obs}}$ & 0.0625 & Global & Observation noise scale \\
$\varepsilon$ (ELBO tol.) & $10^{-8}$ & Global & Standard convergence criterion \\
Max iterations & 100 & Global & Upper bound (converges in 2) \\
$\tau$ (discretization) & 0.5 & Global & Neutral threshold \\
\bottomrule
\end{tabular}
\end{table}

\begin{table}[H]
\centering
\caption{Per-rule prior parameters.}
\label{tab:rule_priors}
\renewcommand{\arraystretch}{1.15}
\begin{tabular}{llccl}
\toprule
\textbf{Rule} & \textbf{Description} & $\alpha_i$ & $\beta_i$
  & $\mathbb{E}[c_i]$ \\
\midrule
R1\_TVA  & Value Added Tax          & 7 & 3 & 0.70 \\
R2\_IS   & Corporate Income Tax     & 6 & 4 & 0.60 \\
R3\_IMF  & Minimum Tax              & 7 & 3 & 0.70 \\
R4\_TPU  & Informal Sector Tax      & 5 & 5 & 0.50 \\
R5\_IRPP & Income Tax Withholding   & 6 & 4 & 0.60 \\
R6\_PAT  & Business License         & 7 & 3 & 0.70 \\
R7\_DECL & Annual Declaration       & 6 & 4 & 0.60 \\
R8\_BANK & Bank Formalization       & 5 & 5 & 0.50 \\
\bottomrule
\end{tabular}
\end{table}

%% ─────────────────────────────────────────────────────────────
%% A.2 — DATASET DOCUMENTATION
%% ─────────────────────────────────────────────────────────────

\subsection{RSI-Togo-Fiscal-Synthetic v2.0}

\paragraph{Overview.}
The benchmark contains 2{,}000 synthetic enterprises across two
regulatory periods (2022--2024 and 2025), encoding 8 fiscal rules.
Each enterprise record is structured in four layers: observable
features (segment, declared turnover), latent ground truth
($c_i^*$ per rule), noisy compliance signals with 18\% missing data,
and binary labels for baseline evaluation only.

\paragraph{Market segments.}
Enterprises are assigned to four segments reflecting the structure
of the Togolese market. Turnover follows a log-normal distribution
within each segment, whose aggregate approximates a power
law~\citep{gabaix2009power}.

\begin{center}
\renewcommand{\arraystretch}{1.15}
\begin{tabular}{lcccc}
\toprule
\textbf{Segment} & \textbf{Share} & \textbf{$\mu_{\log}$}
  & \textbf{$\sigma_{\log}$} & \textbf{Median CA (FCFA)} \\
\midrule
Informal      & 45\% & 16.3 & 0.7 & $\sim$12M \\
Small formal  & 25\% & 17.6 & 0.5 & $\sim$44M \\
Medium        & 18\% & 18.6 & 0.4 & $\sim$119M \\
Large         & 12\% & 20.0 & 0.5 & $\sim$485M \\
\bottomrule
\end{tabular}
\end{center}

\paragraph{Under-declaration model.}
Observed turnover is systematically lower than true turnover:
\[
  \hat{x} = x \cdot \underbrace{\beta}_{\text{systematic}} \cdot
  \underbrace{\varepsilon}_{\text{noise}}, \quad
  \beta \sim \operatorname{Beta}(7, 3), \quad
  \varepsilon \sim \mathcal{N}(1,\, 0.04^2)
\]
The Beta(7,3) distribution has mean $0.70$ and concentrates mass
between 0.50 and 0.90, consistent with empirical estimates from
Sub-Saharan African fiscal studies.

\paragraph{Compliance signals.}
For each entity $j$ and applicable rule $r_i$, the continuous
compliance signal is drawn from
$\operatorname{Beta}(c_{ji} \cdot \kappa,\, (1{-}c_{ji}) \cdot \kappa)$
with $\kappa = 10$, ensuring an unbiased signal centered on the true
entity compliance rate: $\mathbb{E}[s_{ji}^{\mathrm{raw}}] = c_{ji}$.
For R4\_TPU (informal sector), the signal is binary:
$s_{ji} \sim \operatorname{Bernoulli}(c_{ji})$.

\paragraph{Applicability and rules.}
Rules R7\_DECL and R8\_BANK apply universally. The remaining 6 rules
have turnover-based thresholds, ensuring every entity is subject to
at least 3 rules (K\,$\geq$\,3). The distribution is:
K\,=\,3 (56\%), K\,=\,5 (15\%), K\,=\,6 (10\%), K\,=\,7 (19\%).

\begin{center}
\renewcommand{\arraystretch}{1.15}
\begin{tabular}{llll}
\toprule
\textbf{Rule} & \textbf{Description} & \textbf{Condition} & \textbf{Threshold} \\
\midrule
R1\_TVA  & Value Added Tax       & CA $\geq$ threshold & 60M (p1), 100M (p2) \\
R2\_IS   & Corporate Income Tax  & CA $\geq$ threshold & 100M \\
R3\_IMF  & Minimum Tax           & CA $\geq$ threshold & 60M \\
R4\_TPU  & Informal Sector Tax   & CA $<$ threshold    & 60M \\
R5\_IRPP & Income Tax Withholding & CA $\geq$ threshold & 30M \\
R6\_PAT  & Business License      & CA $\geq$ threshold & 30M \\
R7\_DECL & Annual Declaration    & Always              & (universal) \\
R8\_BANK & Bank Formalization    & Always              & (universal) \\
\bottomrule
\end{tabular}
\end{center}

\paragraph{Regulatory change event.}
Period~1 (2022--2024) uses a VAT threshold of 60M~FCFA. Period~2
(2025) uses 100M~FCFA, reflecting Law n\textdegree{}2024-007
(30~December~2024). This provides a clean natural experiment for
validating Theorem~\ref{thm:o1}.

\paragraph{Column structure.}
Each enterprise record contains 48 columns: 6 identifiers and
features (\texttt{entity\_id}, \texttt{period}, \texttt{segment},
\texttt{true\_ca}, \texttt{obs\_ca}, \texttt{under\_decl\_ratio}),
5 columns per rule (applicability, true compliance, raw signal,
discretized signal, label), and 2 aggregate columns
(\texttt{label\_any\_nc}, \texttt{n\_applicable}).

%% ─────────────────────────────────────────────────────────────
%% A.3 — BASELINES
%% ─────────────────────────────────────────────────────────────

\subsection{Baseline Configurations}

\begin{table}[H]
\centering
\caption{Baseline model configurations.}
\label{tab:baselines}
\renewcommand{\arraystretch}{1.15}
\begin{tabular}{lll}
\toprule
\textbf{Model} & \textbf{Parameter} & \textbf{Value} \\
\midrule
LR      & solver & lbfgs \\
        & max\_iter & 1000 \\
        & feature scaling & StandardScaler \\
        & train/test split & 70\%/30\%, stratified \\[4pt]
XGBoost & n\_estimators & 150 \\
        & max\_depth & 4 \\
        & learning\_rate & 0.1 (default) \\
        & train/test split & 70\%/30\%, stratified \\[4pt]
MLP     & hidden\_layer\_sizes & (128, 64, 32) \\
        & activation & ReLU \\
        & max\_iter & 1000 \\
        & early\_stopping & True \\
        & learning\_rate & adaptive \\
        & feature scaling & StandardScaler \\[4pt]
RBS     & Logic & Flag if $s_{ji}^{\mathrm{disc}} = 0$ for any applicable rule \\
        & Missing handling & Ignore (no flag) \\
\bottomrule
\end{tabular}
\end{table}

%% ─────────────────────────────────────────────────────────────
%% A.4 — REPRODUCIBILITY
%% ─────────────────────────────────────────────────────────────

\subsection{Reproducibility}

All experiments use a fixed random seed (42) for full deterministic
reproducibility. The implementation uses Python~3.11 with NumPy,
SciPy, scikit-learn, and pandas. No GPU is required. The full
experiment pipeline completes in under 30~seconds on standard
hardware. The dataset, framework, and all experiment code are
released publicly under the MIT license.

%% ─────────────────────────────────────────────────────────────
%% A.5 — DETAILED PROOFS
%% ─────────────────────────────────────────────────────────────

\subsection{Detailed Proof of Theorem~\ref{thm:o1} ($O(1)$ Adaptability)}

\begin{proof}
Let $\mathcal{U}_k$ be a regulatory update that changes the parameters
of rule $r_k$ from $\Theta_k$ to $\Theta_k'$.

\textit{Step 1: Prior factorization.}
The RSI prior factorizes as
$P(\mathbf{S}) = \prod_{i=1}^n P(c_i \mid \Theta_i) \cdot P(\delta_i \mid \Theta_i)$.
The posterior satisfies:
\[
  P(\mathbf{S} \mid \mathbf{D}) \propto
  \prod_{i=1}^n P(\mathbf{D}_i \mid c_i) \cdot P(c_i \mid \Theta_i) \cdot P(\delta_i \mid \Theta_i)
\]
After the update $\Theta_k \to \Theta_k'$, only the $k$-th factor
changes:
\[
  P'(\mathbf{S} \mid \mathbf{D}) \propto
  P(\mathbf{S} \mid \mathbf{D}) \cdot
  \frac{P(c_k \mid \Theta_k') \cdot P(\delta_k \mid \Theta_k')}
       {P(c_k \mid \Theta_k) \cdot P(\delta_k \mid \Theta_k)}
\]

\textit{Step 2: Drift correction.}
For a threshold change $\theta_k \to \theta_k'$, define
$\varepsilon = (\theta_k' - \theta_k)/\theta_k$. The drift posterior
mean updates as $\bar{\mu}_k \leftarrow \bar{\mu}_k - \varepsilon$.
This is a single scalar subtraction: $O(1)$.

\textit{Step 3: Invariance of compliance parameters.}
The compliance rate $c_k$ is independent of the threshold magnitude
(it measures the \emph{proportion} of applicable entities that comply,
not the threshold itself). Therefore $\bar{\alpha}_k$ and
$\bar{\beta}_k$ are invariant under $\mathcal{U}_k$. Similarly,
$\bar{\sigma}_k$ depends on the number of observations and prior
variance, neither of which changes under $\mathcal{U}_k$.

\textit{Step 4: Normalization.}
Under conjugate families, the posterior normalization constant
$B(\bar{\alpha}_k, \bar{\beta}_k)$ is computed analytically and does
not require re-processing the data. The updated posterior is a
properly normalized Beta distribution.

\textit{Total cost:} one scalar subtraction ($\bar{\mu}_k$), zero
data access. $\mathcal{C}^{\mathrm{RSI}}(\mathcal{U}_k) = O(1)$
with respect to $|\mathbf{D}|$ and $n$. \qed
\end{proof}

\subsection{Detailed Proof of Theorem~\ref{thm:bvm} (BvM Consistency)}

\begin{proof}
We prove the result for a single rule $r_i$ (the extension to
$\mathbf{S}$ follows by factorization).

Let $s_1, s_2, \ldots, s_m$ be i.i.d.\ Bernoulli($c_i^*$) compliance
signals, where $c_i^* \in (0,1)$ is the true population compliance rate.

\textit{Step 1: Posterior computation.}
The conjugate posterior is
$Q^*(c_i) = \operatorname{Beta}(\bar{\alpha}_i, \bar{\beta}_i)$ with
$\bar{\alpha}_i = \alpha_i + \sum_{j=1}^m s_j$ and
$\bar{\beta}_i = \beta_i + \sum_{j=1}^m (1-s_j)$.

\textit{Step 2: Posterior mean convergence.}
Define $\bar{s} = m^{-1}\sum_j s_j$. Then:
\[
  \mathbb{E}_Q[c_i] = \frac{\bar{\alpha}_i}{\bar{\alpha}_i + \bar{\beta}_i}
  = \frac{\alpha_i + m\bar{s}}{\alpha_i + \beta_i + m}
  \to \bar{s} \to c_i^* \quad (m \to \infty)
\]
where the first convergence follows from $m \gg \alpha_i + \beta_i$
and the second from the strong law of large numbers.

\textit{Step 3: Posterior variance.}
\[
  \mathrm{Var}_Q[c_i] = \frac{\bar{\alpha}_i \bar{\beta}_i}
  {(\bar{\alpha}_i + \bar{\beta}_i)^2 (\bar{\alpha}_i + \bar{\beta}_i + 1)}
  \sim \frac{c_i^*(1-c_i^*)}{m} \quad (m \to \infty)
\]
This is the Bernoulli variance divided by $m$, confirming
$\sigma[c_i] \sim 1/\sqrt{m}$.

\textit{Step 4: Asymptotic normality.}
By the central limit theorem, $\sqrt{m}(\bar{s} - c_i^*)
\xrightarrow{d} \mathcal{N}(0, c_i^*(1-c_i^*))$. Since
$\mathbb{E}_Q[c_i] \to \bar{s}$ and $\mathrm{Var}_Q[c_i] \to
c_i^*(1-c_i^*)/m$, we obtain:
\[
  \sqrt{m}(\mathbb{E}_Q[c_i] - c_i^*) \xrightarrow{d}
  \mathcal{N}(0, \mathcal{I}(c_i^*)^{-1})
\]
where $\mathcal{I}(c_i^*) = [c_i^*(1-c_i^*)]^{-1}$ is the Fisher
information of the Bernoulli model. This is the classical
Bernstein-von Mises theorem applied to the Beta-Bernoulli conjugate
pair \citep{van2000asymptotic}. \qed
\end{proof}

%% ─────────────────────────────────────────────────────────────
%% A.6 — SIGNAL DISCRETIZATION BIAS
%% ─────────────────────────────────────────────────────────────

\subsection{Signal Discretization and Calibration Bias}

The compliance signal $s_{ji}^{\mathrm{raw}} \in [0,1]$ is unbiased
by construction:
$\mathbb{E}[s_{ji}^{\mathrm{raw}}] = c_{ji}$.
Discretization at threshold $\tau = 0.5$ introduces a systematic bias.
For an entity with true compliance $c_{ji}$, the probability that the
discretized signal equals 1 is:
\[
  P(s_{ji}^{\mathrm{disc}} = 1) =
  1 - I_{0.5}(c_{ji} \cdot \kappa,\, (1{-}c_{ji}) \cdot \kappa)
\]
where $I_x(a,b)$ is the regularized incomplete Beta function. When
$c_{ji} > 0.5$, this probability exceeds $c_{ji}$ (the discrete signal
overestimates compliance). When $c_{ji} < 0.5$, it underestimates.
The bias is largest for $c_{ji}$ far from 0.5 and decreases as
$\kappa \to \infty$.

This explains the calibration misses in Table~\ref{tab:calibration}:
rules with true population compliance far from 0.5 (R3\_IMF at 0.876,
R7\_DECL at 0.666) exhibit the largest bias, while rules near 0.5
(R4\_TPU at 0.345, R8\_BANK at 0.491) are well-calibrated.

A natural mitigation is to use continuous signals directly with a
Beta-likelihood formulation, eliminating discretization entirely while
preserving conjugacy under moment-matching approximations. This
extension is developed in future work.

%% ─────────────────────────────────────────────────────────────
%% A.7 — PRIOR SENSITIVITY ANALYSIS
%% ─────────────────────────────────────────────────────────────

\subsection{Prior Sensitivity Analysis}

To assess robustness to prior misspecification, we evaluate RSI under
five prior configurations ranging from domain-informed to deliberately
adversarial:

\begin{center}
\renewcommand{\arraystretch}{1.15}
\begin{tabular}{llcc}
\toprule
\textbf{Configuration} & \textbf{Prior} & $\mathbb{E}[c_i]$ & \textbf{Rationale} \\
\midrule
Correct       & $\mathrm{Beta}(7, 3)$ & 0.70 & Domain-informed (OTR) \\
Uniform       & $\mathrm{Beta}(1, 1)$ & 0.50 & Maximum ignorance \\
Inverted      & $\mathrm{Beta}(3, 7)$ & 0.30 & Deliberately wrong direction \\
Strong wrong  & $\mathrm{Beta}(20, 2)$ & 0.91 & Strong confidence, wrong value \\
Weak          & $\mathrm{Beta}(2, 2)$ & 0.50 & Weak, symmetric \\
\bottomrule
\end{tabular}
\end{center}

Entity-level F1 is perfectly invariant (0.741 across all
configurations) because entity scoring is deterministic:
$\mathrm{NC}_{ji} = 1 - s_{ji}^{\mathrm{raw}}$ does not depend on
the prior. Population posteriors $\mathbb{E}[c_i]$ vary by at most
0.027 across configurations, confirming that with $n > 400$
observations per rule, the likelihood dominates the prior. This
validates the prior robustness corollary of Theorem~\ref{thm:bvm}:
different initial beliefs converge to the same assessment as evidence
accumulates.
\section{Appendix: Supplementary Material}
\label{sec:appendix}

%% ─────────────────────────────────────────────────────────────
%% A.1 — HYPERPARAMETERS
%% ─────────────────────────────────────────────────────────────

\subsection{Hyperparameters}

Table~\ref{tab:hyperparams} lists all hyperparameters used in the RSI
Togolese fiscal instantiation. All values are derived from institutional
knowledge (OTR fiscal rules, historical compliance estimates) or set to
standard non-informative values. No hyperparameter is learned from
labeled data.

\begin{table}[H]
\centering
\caption{RSI hyperparameters. Prior parameters encode institutional
         knowledge; inference parameters are standard.}
\label{tab:hyperparams}
\renewcommand{\arraystretch}{1.15}
\begin{tabular}{llll}
\toprule
\textbf{Parameter} & \textbf{Value} & \textbf{Scope} & \textbf{Justification} \\
\midrule
\multicolumn{4}{l}{\textit{Prior hyperparameters (per rule)}} \\
$\alpha_i, \beta_i$ & See Table~\ref{tab:rule_priors} & Per rule & OTR historical rates \\
$\sigma_i$ (drift) & 0.5--2.0 & Per rule & Expected parametric stability \\
\midrule
\multicolumn{4}{l}{\textit{Inference parameters}} \\
$\sigma_{\mathrm{obs}}$ & 0.0625 & Global & Observation noise scale \\
$\varepsilon$ (ELBO tol.) & $10^{-8}$ & Global & Standard convergence criterion \\
Max iterations & 100 & Global & Upper bound (converges in 2) \\
$\tau$ (discretization) & 0.5 & Global & Neutral threshold \\
\bottomrule
\end{tabular}
\end{table}

\begin{table}[H]
\centering
\caption{Per-rule prior parameters.}
\label{tab:rule_priors}
\renewcommand{\arraystretch}{1.15}
\begin{tabular}{llccl}
\toprule
\textbf{Rule} & \textbf{Description} & $\alpha_i$ & $\beta_i$
  & $\mathbb{E}[c_i]$ \\
\midrule
R1\_TVA  & Value Added Tax          & 7 & 3 & 0.70 \\
R2\_IS   & Corporate Income Tax     & 6 & 4 & 0.60 \\
R3\_IMF  & Minimum Tax              & 7 & 3 & 0.70 \\
R4\_TPU  & Informal Sector Tax      & 5 & 5 & 0.50 \\
R5\_IRPP & Income Tax Withholding   & 6 & 4 & 0.60 \\
R6\_PAT  & Business License         & 7 & 3 & 0.70 \\
R7\_DECL & Annual Declaration       & 6 & 4 & 0.60 \\
R8\_BANK & Bank Formalization       & 5 & 5 & 0.50 \\
\bottomrule
\end{tabular}
\end{table}

%% ─────────────────────────────────────────────────────────────
%% A.2 — DATASET DOCUMENTATION
%% ─────────────────────────────────────────────────────────────

\subsection{RSI-Togo-Fiscal-Synthetic v2.0}

\paragraph{Overview.}
The benchmark contains 2{,}000 synthetic enterprises across two
regulatory periods (2022--2024 and 2025), encoding 8 fiscal rules.
Each enterprise record is structured in four layers: observable
features (segment, declared turnover), latent ground truth
($c_i^*$ per rule), noisy compliance signals with 18\% missing data,
and binary labels for baseline evaluation only.

\paragraph{Market segments.}
Enterprises are assigned to four segments reflecting the structure
of the Togolese market. Turnover follows a log-normal distribution
within each segment, whose aggregate approximates a power
law~\citep{gabaix2009power}.

\begin{center}
\renewcommand{\arraystretch}{1.15}
\begin{tabular}{lcccc}
\toprule
\textbf{Segment} & \textbf{Share} & \textbf{$\mu_{\log}$}
  & \textbf{$\sigma_{\log}$} & \textbf{Median CA (FCFA)} \\
\midrule
Informal      & 45\% & 16.3 & 0.7 & $\sim$12M \\
Small formal  & 25\% & 17.6 & 0.5 & $\sim$44M \\
Medium        & 18\% & 18.6 & 0.4 & $\sim$119M \\
Large         & 12\% & 20.0 & 0.5 & $\sim$485M \\
\bottomrule
\end{tabular}
\end{center}

\paragraph{Under-declaration model.}
Observed turnover is systematically lower than true turnover:
\[
  \hat{x} = x \cdot \underbrace{\beta}_{\text{systematic}} \cdot
  \underbrace{\varepsilon}_{\text{noise}}, \quad
  \beta \sim \operatorname{Beta}(7, 3), \quad
  \varepsilon \sim \mathcal{N}(1,\, 0.04^2)
\]
The Beta(7,3) distribution has mean $0.70$ and concentrates mass
between 0.50 and 0.90, consistent with empirical estimates from
Sub-Saharan African fiscal studies.

\paragraph{Compliance signals.}
For each entity $j$ and applicable rule $r_i$, the continuous
compliance signal is drawn from
$\operatorname{Beta}(c_{ji} \cdot \kappa,\, (1{-}c_{ji}) \cdot \kappa)$
with $\kappa = 10$, ensuring an unbiased signal centered on the true
entity compliance rate: $\mathbb{E}[s_{ji}^{\mathrm{raw}}] = c_{ji}$.
For R4\_TPU (informal sector), the signal is binary:
$s_{ji} \sim \operatorname{Bernoulli}(c_{ji})$.

\paragraph{Applicability and rules.}
Rules R7\_DECL and R8\_BANK apply universally. The remaining 6 rules
have turnover-based thresholds, ensuring every entity is subject to
at least 3 rules (K\,$\geq$\,3). The distribution is:
K\,=\,3 (56\%), K\,=\,5 (15\%), K\,=\,6 (10\%), K\,=\,7 (19\%).

\begin{center}
\renewcommand{\arraystretch}{1.15}
\begin{tabular}{llll}
\toprule
\textbf{Rule} & \textbf{Description} & \textbf{Condition} & \textbf{Threshold} \\
\midrule
R1\_TVA  & Value Added Tax       & CA $\geq$ threshold & 60M (p1), 100M (p2) \\
R2\_IS   & Corporate Income Tax  & CA $\geq$ threshold & 100M \\
R3\_IMF  & Minimum Tax           & CA $\geq$ threshold & 60M \\
R4\_TPU  & Informal Sector Tax   & CA $<$ threshold    & 60M \\
R5\_IRPP & Income Tax Withholding & CA $\geq$ threshold & 30M \\
R6\_PAT  & Business License      & CA $\geq$ threshold & 30M \\
R7\_DECL & Annual Declaration    & Always              & (universal) \\
R8\_BANK & Bank Formalization    & Always              & (universal) \\
\bottomrule
\end{tabular}
\end{center}

\paragraph{Regulatory change event.}
Period~1 (2022--2024) uses a VAT threshold of 60M~FCFA. Period~2
(2025) uses 100M~FCFA, reflecting Law n\textdegree{}2024-007
(30~December~2024). This provides a clean natural experiment for
validating Theorem~\ref{thm:o1}.

\paragraph{Column structure.}
Each enterprise record contains 48 columns: 6 identifiers and
features (\texttt{entity\_id}, \texttt{period}, \texttt{segment},
\texttt{true\_ca}, \texttt{obs\_ca}, \texttt{under\_decl\_ratio}),
5 columns per rule (applicability, true compliance, raw signal,
discretized signal, label), and 2 aggregate columns
(\texttt{label\_any\_nc}, \texttt{n\_applicable}).

%% ─────────────────────────────────────────────────────────────
%% A.3 — BASELINES
%% ─────────────────────────────────────────────────────────────

\subsection{Baseline Configurations}

\begin{table}[H]
\centering
\caption{Baseline model configurations.}
\label{tab:baselines}
\renewcommand{\arraystretch}{1.15}
\begin{tabular}{lll}
\toprule
\textbf{Model} & \textbf{Parameter} & \textbf{Value} \\
\midrule
LR      & solver & lbfgs \\
        & max\_iter & 1000 \\
        & feature scaling & StandardScaler \\
        & train/test split & 70\%/30\%, stratified \\[4pt]
XGBoost & n\_estimators & 150 \\
        & max\_depth & 4 \\
        & learning\_rate & 0.1 (default) \\
        & train/test split & 70\%/30\%, stratified \\[4pt]
MLP     & hidden\_layer\_sizes & (128, 64, 32) \\
        & activation & ReLU \\
        & max\_iter & 1000 \\
        & early\_stopping & True \\
        & learning\_rate & adaptive \\
        & feature scaling & StandardScaler \\[4pt]
RBS     & Logic & Flag if $s_{ji}^{\mathrm{disc}} = 0$ for any applicable rule \\
        & Missing handling & Ignore (no flag) \\
\bottomrule
\end{tabular}
\end{table}

%% ─────────────────────────────────────────────────────────────
%% A.4 — REPRODUCIBILITY
%% ─────────────────────────────────────────────────────────────

\subsection{Reproducibility}

All experiments use a fixed random seed (42) for full deterministic
reproducibility. The implementation uses Python~3.11 with NumPy,
SciPy, scikit-learn, and pandas. No GPU is required. The full
experiment pipeline completes in under 30~seconds on standard
hardware. The dataset, framework, and all experiment code are
released publicly under the MIT license.

%% ─────────────────────────────────────────────────────────────
%% A.5 — DETAILED PROOFS
%% ─────────────────────────────────────────────────────────────

\subsection{Detailed Proof of Theorem~\ref{thm:o1} ($O(1)$ Adaptability)}

\begin{proof}
Let $\mathcal{U}_k$ be a regulatory update that changes the parameters
of rule $r_k$ from $\Theta_k$ to $\Theta_k'$.

\textit{Step 1: Prior factorization.}
The RSI prior factorizes as
$P(\mathbf{S}) = \prod_{i=1}^n P(c_i \mid \Theta_i) \cdot P(\delta_i \mid \Theta_i)$.
The posterior satisfies:
\[
  P(\mathbf{S} \mid \mathbf{D}) \propto
  \prod_{i=1}^n P(\mathbf{D}_i \mid c_i) \cdot P(c_i \mid \Theta_i) \cdot P(\delta_i \mid \Theta_i)
\]
After the update $\Theta_k \to \Theta_k'$, only the $k$-th factor
changes:
\[
  P'(\mathbf{S} \mid \mathbf{D}) \propto
  P(\mathbf{S} \mid \mathbf{D}) \cdot
  \frac{P(c_k \mid \Theta_k') \cdot P(\delta_k \mid \Theta_k')}
       {P(c_k \mid \Theta_k) \cdot P(\delta_k \mid \Theta_k)}
\]

\textit{Step 2: Drift correction.}
For a threshold change $\theta_k \to \theta_k'$, define
$\varepsilon = (\theta_k' - \theta_k)/\theta_k$. The drift posterior
mean updates as $\bar{\mu}_k \leftarrow \bar{\mu}_k - \varepsilon$.
This is a single scalar subtraction: $O(1)$.

\textit{Step 3: Invariance of compliance parameters.}
The compliance rate $c_k$ is independent of the threshold magnitude
(it measures the \emph{proportion} of applicable entities that comply,
not the threshold itself). Therefore $\bar{\alpha}_k$ and
$\bar{\beta}_k$ are invariant under $\mathcal{U}_k$. Similarly,
$\bar{\sigma}_k$ depends on the number of observations and prior
variance, neither of which changes under $\mathcal{U}_k$.

\textit{Step 4: Normalization.}
Under conjugate families, the posterior normalization constant
$B(\bar{\alpha}_k, \bar{\beta}_k)$ is computed analytically and does
not require re-processing the data. The updated posterior is a
properly normalized Beta distribution.

\textit{Total cost:} one scalar subtraction ($\bar{\mu}_k$), zero
data access. $\mathcal{C}^{\mathrm{RSI}}(\mathcal{U}_k) = O(1)$
with respect to $|\mathbf{D}|$ and $n$. \qed
\end{proof}

\subsection{Detailed Proof of Theorem~\ref{thm:bvm} (BvM Consistency)}

\begin{proof}
We prove the result for a single rule $r_i$ (the extension to
$\mathbf{S}$ follows by factorization).

Let $s_1, s_2, \ldots, s_m$ be i.i.d.\ Bernoulli($c_i^*$) compliance
signals, where $c_i^* \in (0,1)$ is the true population compliance rate.

\textit{Step 1: Posterior computation.}
The conjugate posterior is
$Q^*(c_i) = \operatorname{Beta}(\bar{\alpha}_i, \bar{\beta}_i)$ with
$\bar{\alpha}_i = \alpha_i + \sum_{j=1}^m s_j$ and
$\bar{\beta}_i = \beta_i + \sum_{j=1}^m (1-s_j)$.

\textit{Step 2: Posterior mean convergence.}
Define $\bar{s} = m^{-1}\sum_j s_j$. Then:
\[
  \mathbb{E}_Q[c_i] = \frac{\bar{\alpha}_i}{\bar{\alpha}_i + \bar{\beta}_i}
  = \frac{\alpha_i + m\bar{s}}{\alpha_i + \beta_i + m}
  \to \bar{s} \to c_i^* \quad (m \to \infty)
\]
where the first convergence follows from $m \gg \alpha_i + \beta_i$
and the second from the strong law of large numbers.

\textit{Step 3: Posterior variance.}
\[
  \mathrm{Var}_Q[c_i] = \frac{\bar{\alpha}_i \bar{\beta}_i}
  {(\bar{\alpha}_i + \bar{\beta}_i)^2 (\bar{\alpha}_i + \bar{\beta}_i + 1)}
  \sim \frac{c_i^*(1-c_i^*)}{m} \quad (m \to \infty)
\]
This is the Bernoulli variance divided by $m$, confirming
$\sigma[c_i] \sim 1/\sqrt{m}$.

\textit{Step 4: Asymptotic normality.}
By the central limit theorem, $\sqrt{m}(\bar{s} - c_i^*)
\xrightarrow{d} \mathcal{N}(0, c_i^*(1-c_i^*))$. Since
$\mathbb{E}_Q[c_i] \to \bar{s}$ and $\mathrm{Var}_Q[c_i] \to
c_i^*(1-c_i^*)/m$, we obtain:
\[
  \sqrt{m}(\mathbb{E}_Q[c_i] - c_i^*) \xrightarrow{d}
  \mathcal{N}(0, \mathcal{I}(c_i^*)^{-1})
\]
where $\mathcal{I}(c_i^*) = [c_i^*(1-c_i^*)]^{-1}$ is the Fisher
information of the Bernoulli model. This is the classical
Bernstein-von Mises theorem applied to the Beta-Bernoulli conjugate
pair \citep{van2000asymptotic}. \qed
\end{proof}

%% ─────────────────────────────────────────────────────────────
%% A.6 — SIGNAL DISCRETIZATION BIAS
%% ─────────────────────────────────────────────────────────────

\subsection{Signal Discretization and Calibration Bias}

The compliance signal $s_{ji}^{\mathrm{raw}} \in [0,1]$ is unbiased
by construction:
$\mathbb{E}[s_{ji}^{\mathrm{raw}}] = c_{ji}$.
Discretization at threshold $\tau = 0.5$ introduces a systematic bias.
For an entity with true compliance $c_{ji}$, the probability that the
discretized signal equals 1 is:
\[
  P(s_{ji}^{\mathrm{disc}} = 1) =
  1 - I_{0.5}(c_{ji} \cdot \kappa,\, (1{-}c_{ji}) \cdot \kappa)
\]
where $I_x(a,b)$ is the regularized incomplete Beta function. When
$c_{ji} > 0.5$, this probability exceeds $c_{ji}$ (the discrete signal
overestimates compliance). When $c_{ji} < 0.5$, it underestimates.
The bias is largest for $c_{ji}$ far from 0.5 and decreases as
$\kappa \to \infty$.

This explains the calibration misses in Table~\ref{tab:calibration}:
rules with true population compliance far from 0.5 (R3\_IMF at 0.876,
R7\_DECL at 0.666) exhibit the largest bias, while rules near 0.5
(R4\_TPU at 0.345, R8\_BANK at 0.491) are well-calibrated.

A natural mitigation is to use continuous signals directly with a
Beta-likelihood formulation, eliminating discretization entirely while
preserving conjugacy under moment-matching approximations. This
extension is developed in future work.

%% ─────────────────────────────────────────────────────────────
%% A.7 — PRIOR SENSITIVITY ANALYSIS
%% ─────────────────────────────────────────────────────────────

\subsection{Prior Sensitivity Analysis}

To assess robustness to prior misspecification, we evaluate RSI under
five prior configurations ranging from domain-informed to deliberately
adversarial:

\begin{center}
\renewcommand{\arraystretch}{1.15}
\begin{tabular}{llcc}
\toprule
\textbf{Configuration} & \textbf{Prior} & $\mathbb{E}[c_i]$ & \textbf{Rationale} \\
\midrule
Correct       & $\mathrm{Beta}(7, 3)$ & 0.70 & Domain-informed (OTR) \\
Uniform       & $\mathrm{Beta}(1, 1)$ & 0.50 & Maximum ignorance \\
Inverted      & $\mathrm{Beta}(3, 7)$ & 0.30 & Deliberately wrong direction \\
Strong wrong  & $\mathrm{Beta}(20, 2)$ & 0.91 & Strong confidence, wrong value \\
Weak          & $\mathrm{Beta}(2, 2)$ & 0.50 & Weak, symmetric \\
\bottomrule
\end{tabular}
\end{center}

Entity-level F1 is perfectly invariant (0.741 across all
configurations) because entity scoring is deterministic:
$\mathrm{NC}_{ji} = 1 - s_{ji}^{\mathrm{raw}}$ does not depend on
the prior. Population posteriors $\mathbb{E}[c_i]$ vary by at most
0.027 across configurations, confirming that with $n > 400$
observations per rule, the likelihood dominates the prior. This
validates the prior robustness corollary of Theorem~\ref{thm:bvm}:
different initial beliefs converge to the same assessment as evidence
accumulates.

\end{document}